\begin{solapartat}
\punts{0,5} Sabem que la imatge és de la mateixa mida però està invertida, llavors:
\[ \beta = \frac{y'}{y} = -1 \]
A més:
\[ \beta = \frac{y'}{y} = \frac{s'}{s} = -1 \Longrightarrow s' = -s \]
I també sabem que $s = -20\ \mathrm{cm}$, llavors $s' = 20\ \mathrm{cm}$, per tant, cal col·locar el sensor a 20~cm de la lent.

\punts{0,5} Per trobar la focal de la lent, farem servir l'equació de la lent:
\[ \frac{1}{f'} = \frac{1}{s'} - \frac{1}{s} = \frac{1}{20} - \frac{1}{-20} = \frac{1}{10} \Longrightarrow f' = 10\ \mathrm{cm} \]
\punts{0,25} Com que la longitud focal imatge és positiva, llavors la lent ha de ser \textbf{convergent}. També es pot justificar que la focal imatge està situada a la dreta de la lent i, per tant, que la lent és convergent.
\end{solapartat}

\begin{solapartat}
\punts{1,25} Per al diagrama de rajos només cal dibuixar 2 dels 3 rajos indicats a la figura: el raig que surt de l'objecte i és paral·lel a l'eix passa pel punt focal imatge F', el raig que passa pel centre de la lent i no es desvia, i el raig que surt de l'objecte passa per F i surt paral·lel a l'eix.

\figura[0.9]{sol_fig1.png}

\destaca{També es considerarà correcte que es representi l'objecte invertit i la imatge dreta.}
\end{solapartat}
